
\documentclass[10pt,landscape]{article}
\usepackage{mathwizards-formulario}

% ── Comandos propios de este documento ─────────
\renewcommand{\headrulewidth}{0pt}
\pagestyle{fancy}
\fancyhf{}
\fancyfoot[C]{\color{primary}\small\textbf{F\'{i}sica I -- Trabajo, Energ\'{i}a y Calor | UDO N\'{u}cleo Monagas}}

\begin{document}
\pagecolor{white}

% Título principal
\begin{center}
{\Huge\bfseries\color{primary} FORMULARIO DE TRABAJO, ENERG\'{I}A Y CALOR}\\[2pt]
{\large\color{darkbg}\textbf{Gu\'{i}a de Estudio V} | Hoja Carta Horizontal}
\end{center}
\vspace{-4pt}

\begin{multicols}{3}

% ======================== COLUMNA 1 ========================

\cajatitulo{1. TRABAJO MEC\'{A}NICO ($W$)}

\subsection*{Trabajo de una Fuerza Constante con \'{A}ngulo}
\[
\boxed{W = F \cdot d \cdot \cos(\theta)}
\]
\begin{itemize}[leftmargin=*,nosep,topsep=0pt]
  \item $F$: magnitud de la fuerza (N)
  \item $d$: desplazamiento (m)
  \item $\theta$: \'{a}ngulo entre fuerza y desplazamiento
\end{itemize}

\textbf{Despejes:}
\[
F = \frac{W}{d \cdot \cos\theta} \quad d = \frac{W}{F \cdot \cos\theta} \quad \theta = \arccos\!\left(\frac{W}{F \cdot d}\right)
\]
\cajainfo{Aplicable a: Ej. 1, 25, 39, 44}

\subsection*{Trabajo de la Fuerza de Gravedad (Peso)}
\[
\boxed{W_g = F_g \cdot d \cdot \cos(\theta_g) = m \cdot g \cdot h}
\]
Si el cuerpo \textbf{baja}: $W_g > 0$ (positivo). Si \textbf{sube}: $W_g < 0$.
\cajainfo{Aplicable a: Ej. 4, 5, 10, 14, 24}

\vspace{4pt}
\cajatitulo{2. POTENCIA MEC\'{A}NICA ($P$)}

\subsection*{Potencia Media}
\[
\boxed{P = \frac{W}{t}}
\]
\textbf{Despejes:}
\[
W = P \cdot t \qquad t = \frac{W}{P}
\]
\cajainfo{Aplicable a: Ej. 4, 5, 10, 14, 15, 17, 18}

\subsection*{Potencia Instant\'{a}nea}
\[
\boxed{P = F \cdot v}
\]
\textbf{Despejes:}
\[
F = \frac{P}{v} \qquad v = \frac{P}{F}
\]
\cajainfo{Aplicable a: Ej. 14, 43}

\vspace{4pt}
\cajatitulo{3. ENERG\'{I}AS MEC\'{A}NICAS}

\subsection*{Energ\'{i}a Cin\'{e}tica ($E_c$)}
\[
\boxed{E_c = \frac{1}{2} m \cdot v^2}
\]
\textbf{Despejes:}
\[
m = \frac{2 \cdot E_c}{v^2} \qquad v = \sqrt{\frac{2 \cdot E_c}{m}}
\]
\cajainfo{Aplicable a: Ej. 7, 16, 18, 21, 22, 27, 33, 38, 42}

\subsection*{Energ\'{i}a Potencial Gravitatoria ($E_p$)}
\[
\boxed{E_p = m \cdot g \cdot h}
\]
\textbf{Despejes:}
\[
m = \frac{E_p}{g \cdot h} \qquad h = \frac{E_p}{m \cdot g}
\]
\cajainfo{Aplicable a: Ej. 11, 19, 22, 33, 36, 38}

\subsection*{Energ\'{i}a Potencial El\'{a}stica ($E_{pe}$)}
\[
\boxed{E_{pe} = \frac{1}{2} k \cdot x^2}
\]
\begin{itemize}[leftmargin=*,nosep,topsep=0pt]
  \item $k$: constante el\'{a}stica (N/m)
  \item $x$: deformaci\'{o}n en metros (m)
\end{itemize}
\textbf{Despejes:}
\[
k = \frac{2 \cdot E_{pe}}{x^2} \qquad x = \sqrt{\frac{2 \cdot E_{pe}}{k}}
\]
\cajainfo{Aplicable a: Ej. 6, 7, 13, 26, 27, 36, 45}

\subsection*{Fuerza El\'{a}stica (Ley de Hooke)}
\[
\boxed{F_e = k \cdot x}
\]
\cajainfo{Util para hallar $k$ en el Ej. 26 antes de calcular la energ\'{i}a}

% ======================== COLUMNA 2 ========================

\cajatitulo{4. TEOREMAS Y CONSERVACI\'{O}N}

\subsection*{Teorema Trabajo -- Energ\'{i}a Cin\'{e}tica}
\[
\boxed{W_{\text{neto}} = \Delta E_c = E_{cf} - E_{ci}}
\]
\cajainfo{Aplicable a: Ej. 1, 3, 16, 18}

\subsection*{Conservaci\'{o}n de la Energ\'{i}a Mec\'{a}nica}
Si solo act\'{u}an fuerzas conservativas:
\[
\boxed{E_{mi} = E_{mf}}
\]
\[
E_{ci} + E_{pi} + E_{pei} = E_{cf} + E_{pf} + E_{pef}
\]
\cajainfo{Aplicable a: Ej. 7, 13, 22, 33, 38}

\subsection*{Teorema General (Con Rozamiento)}
\[
\boxed{W_{nc} = \Delta E_m = E_{mf} - E_{mi}}
\]
Si hay disipaci\'{o}n expl\'{i}cita (ej. ``consume 25\%''):
\[
E_{mf} = E_{mi} - E_{\text{disipada}}
\]
\cajainfo{Aplicable a: Ej. 19, 21, 31, 34, 36, 40, 41, 42, 45}

\subsection*{Choques Inel\'{a}sticos (Momento Lineal)}
Antes de usar conservaci\'{o}n de energ\'{i}a con resortes tras un impacto:
\[
\boxed{m_1 \cdot v_1 = (m_1 + m_2) \cdot v_f}
\]
\cajainfo{Aplicable a: Ej. 27 (bala incrustada en bloque)}

\vspace{4pt}
\cajatitulo{5. DIN\'{A}MICA AUXILIAR}

\subsection*{Segunda Ley de Newton}
\[
\boxed{F_{\text{neta}} = m \cdot a} \qquad \boxed{a = \frac{F_{\text{neta}}}{m}}
\]
\cajainfo{Necesaria para Ej. 3, 15, 28, 35 donde se dan fuerzas y tiempos}

\subsection*{Cinem\'{a}tica B\'{a}sica}
\[
v_f = v_0 + a \cdot t
\]
\[
d = v_0 \cdot t + \frac{1}{2} a \cdot t^2
\]
\[
\boxed{v_f^2 = v_0^2 + 2 \cdot a \cdot d}
\]

\subsection*{Fricci\'{o}n}
\[
f_k = \mu_k \cdot N \qquad f_s \leq \mu_s \cdot N
\]
Trabajo de la fricci\'{o}n (dissipativo):
\[
W_{\text{fricci\'{o}n}} = -f_k \cdot d = -\mu_k \cdot N \cdot d
\]

\vspace{4pt}
\cajatitulo{6. CONVERSIONES DE UNIDADES}

\begin{center}
\begin{tabular}{@{}ll@{}}
\toprule
\textbf{Conversi\'{o}n} & \textbf{F\'{o}rmula} \\
\midrule
Temperatura (°F a °C) & $T(°C) = \frac{5}{9}(T(°F) - 32)$ \\
Volumen a Masa (Agua) & $1\,\text{L} = 1\,\text{kg}$ ($\rho = 1000\,\text{kg/m}^3$) \\
Volumen a Masa (Gasolina) & $\text{Masa} = V \cdot \rho$ \\
Longitud & $1\,\text{m} = 100\,\text{cm} = 1000\,\text{mm}$ \\
Velocidad & $1\,\text{m/s} = 3.6\,\text{km/h}$ \\
Energ\'{i}a & $1\,\text{kcal} = 1000\,\text{cal} \approx 4184\,\text{J}$ \\
\bottomrule
\end{tabular}
\end{center}

\cajainfo{¡Cuidado! La constante $k$ en N/mm (Ej. 27) debe pasarse a N/m}

% ======================== COLUMNA 3 ========================

\cajatitulo{7. RESUMEN DE ESTRATEGIAS}

\subsection*{Toboganes y Monta\~{n}as Rusas (Ej. 21, 22, 33)}
\begin{enumerate}[leftmargin=*,nosep,topsep=0pt]
  \item Elegir nivel de referencia ($h=0$): salida, suelo o punto m\'{a}s bajo.
  \item Aplicar $E_{mi} = E_{mf}$ si no hay rozamiento.
  \item Si hay rozamiento: $W_{\text{fricci\'{o}n}} = E_{mf} - E_{mi}$.
  \item La velocidad final es la misma sin importar d\'{o}nde pongas $h=0$.
\end{enumerate}

\subsection*{Problemas con Resortes (Ej. 6, 7, 13, 26, 27, 36, 45)}
\begin{enumerate}[leftmargin=*,nosep,topsep=0pt]
  \item Si se comprime/estira con fuerza conocida: usar $F_e = kx$ para hallar $k$.
  \item Luego aplicar $E_{pe} = \frac{1}{2}kx^2$.
  \item En choques con resorte: primero conservaci\'{o}n de momento, luego energ\'{i}a.
\end{enumerate}

\subsection*{Problemas con P\'{e}rdidas de Energ\'{i}a (Ej. 19, 21, 31, 34, 36, 40--42, 45)}
\begin{enumerate}[leftmargin=*,nosep,topsep=0pt]
  \item Calcular $E_{mi}$ completa (cin\'{e}tica + potencial + el\'{a}stica).
  \item Aplicar el porcentaje de p\'{e}rdida: $E_{\text{disipada}} = \% \times E_{mi}$.
  \item $E_{mf} = E_{mi} - E_{\text{disipada}}$.
  \item Despejar la inc\'{o}gnita (velocidad, altura, deformaci\'{o}n, etc.).
\end{enumerate}

\subsection*{Problemas de Potencia (Ej. 4, 5, 10, 14--18, 43)}
\begin{enumerate}[leftmargin=*,nosep,topsep=0pt]
  \item Identificar si piden potencia media ($P = W/t$) o instant\'{a}nea ($P = Fv$).
  \item Calcular el trabajo primero si es necesario.
  \item Verificar unidades: W en J, t en s, P en W (J/s).
\end{enumerate}

\vspace{4pt}
\cajatitulo{8. CONSTANTES \'{U}TILES}

\begin{center}
\begin{tabular}{@{}ll@{}}
\toprule
\textbf{Constante} & \textbf{Valor} \\
\midrule
$g$ (gravedad) & $9.8\,\text{m/s}^2$ \\
$G$ (gravitacional) & $6.67 \times 10^{-11}\,\text{N}\cdot\text{m}^2/\text{kg}^2$ \\
Masa Tierra & $5.97 \times 10^{24}\,\text{kg}$ \\
Radio Tierra & $6.37 \times 10^6\,\text{m}$ \\
Calor espec\'{i}fico agua & $4186\,\text{J/kg}\cdot°\text{C}$ \\
Calor latente fusi\'{o}n hielo & $3.34 \times 10^5\,\text{J/kg}$ \\
Calor latente vaporizaci\'{o}n & $2.26 \times 10^6\,\text{J/kg}$ \\
\bottomrule
\end{tabular}
\end{center}

\vspace{4pt}
\cajatitulo{9. CONSEJOS PARA EL EXAMEN}

\begin{enumerate}[leftmargin=*,nosep,topsep=0pt]
  \item \textbf{DCL siempre:} Dibujar el diagrama de cuerpo libre.
  \item \textbf{Nivel de referencia:} Elige $h=0$ donde m\'{a}s te convenga; s\'{e} consistente.
  \item \textbf{Signos de altura:} Por encima de $h=0$ es positiva; por debajo, negativa.
  \item \textbf{Unidades:} Trabajar en SI: kg, m, s, N, J, W.
  \item \textbf{Conversiones:} g $\to$ kg, cm $\to$ m, km/h $\to$ m/s, kcal $\to$ J.
  \item \textbf{Resorte:} La deformaci\'{o}n $x$ siempre en metros, nunca en cm o mm.
  \item \textbf{Rozamiento:} Trabajo negativo (disipa energ\'{i}a). $W_{\text{fricci\'{o}n}} = -f_k \cdot d$.
  \item \textbf{Choques:} Primero momento lineal, despu\'{e}s energ\'{i}a (si aplica).
  \item \textbf{Potencia:} Distinguir entre media e instant\'{a}nea.
\end{enumerate}

\end{multicols}

\end{document}
